%% ma: L10-B13-0023
%% lop: 10
%% chuong: 5
%% bai: 13
%% dang: 10.13.4
%% loai: TN4
%% muc: TH
%% dap_an: "C"
%% nguon: "chuyen-de-cac-so-dac-trung-cua-mau-so-lieu-khong-ghep-nhom-toan-10-chuong-trinh-moi.pdf (GV. Phan Nhật Linh) (P02), Bài 02 Dạng 1 Phần I Câu 13"
%% kien_thuc: "Lớp 10 Bài 13 (trung vị)"
%% trang_thai: chinh-thuc
%% ghi_chu: "Tự duyệt theo yêu cầu giáo viên 22/10/2026."
\begin{ex}
Bạn Danh cân lần lượt $50$ quả vải thiều được chọn ngẫu nhiên từ vườn nhà mình và được kết quả như sau:
\begin{center}
\begin{tabular}{|c|c|c|c|c|c|}\hline
Cân nặng (g) & $8$ & $19$ & $20$ & $21$ & $22$\\\hline
Số quả & $1$ & $10$ & $19$ & $17$ & $3$\\\hline
\end{tabular}
\end{center}
Trung vị của mẫu số liệu trên là
\choice{$19$}{$19{,}5$}{\True $20$}{$21$}
\loigiai{$n=50$ nên $M_e$ là trung bình cộng của giá trị thứ $25$ và thứ $26$. Tần số tích lũy: $1;11;30$ nên cả hai giá trị đều là $20$. Vậy $M_e=20$. Chọn C.}
\end{ex}
